% from ega3-1-fr.tex
\oldpage[0\textsubscript{III}]{66}
tel que
$y_\beta=g_{\beta\lambda}(y_\lambda)+u_\beta(x_\beta)$. On en déduit
\[
y_\alpha=g_{\alpha\lambda}(y_\lambda)
+u_\alpha(f_{\alpha\beta}(x_\beta))~;
\]
mais comme $\lambda\geq\beta$, il existe $x_\lambda\in A_\lambda$ tel que
$f_{\alpha\beta}(x_\beta)=f_{\alpha\lambda}(x_\lambda)$, et finalement
\[
y_\alpha
=g_{\alpha\lambda}(y_\lambda)+u_\alpha(f_{\alpha\lambda}(x_\lambda))
=g_{\alpha\lambda}(y_\lambda+u_\lambda(x_\lambda))
\in g_{\alpha\lambda}(B_\lambda),
\]
ce qui achève la démonstration.
\end{proof}

\begin{proposition}[13.2.2]
\label{0.13.2.2-fr}
\emph{Soit $I$ un ensemble ordonné filtrant ayant une partie cofinale
dénombrable. Soit}
\[
0\longrightarrow A_\alpha\xrightarrow{u_\alpha}B_\alpha
\xrightarrow{v_\alpha}C_\alpha\longrightarrow0
\]
\emph{une suite exacte de systèmes projectifs de groupes abéliens ayant
$I$ pour ensemble d'indices. Si $(A_\alpha)$ satisfait à la condition
(ML), la suite}
\[
0\longrightarrow\varprojlim A_\alpha
\longrightarrow\varprojlim B_\alpha
\longrightarrow\varprojlim C_\alpha\longrightarrow0
\]
\emph{est exacte.}
\end{proposition}

\begin{proof}
Tout revient à prouver que l'homomorphisme
\[
v=\varprojlim v_\alpha:\varprojlim B_\alpha\longrightarrow
\varprojlim C_\alpha
\]
est surjectif. Soit $z=(z_\alpha)$ un élément de
$\varprojlim C_\alpha$, et posons
$E_\alpha=v_\alpha^{-1}(z_\alpha)$~; il est clair que les $E_\alpha$
forment un système projectif d'ensembles non vides pour les restrictions
des homomorphismes $g_{\alpha\beta}:B_\beta\to B_\alpha$. Montrons que ce
système projectif vérifie la condition (ML)~; identifiant $A_\alpha$ à une
partie de $B_\alpha$ par $u_\alpha$, pour tout $\alpha\in I$, il existe
$\beta\geq\alpha$ tel que
$g_{\alpha\beta}(A_\beta)=g_{\alpha\lambda}(A_\lambda)$ pour
$\lambda\geq\beta$~; montrons que l'on a aussi
$g_{\alpha\beta}(E_\beta)=g_{\alpha\lambda}(E_\lambda)$ pour
$\lambda\geq\beta$. En effet, prenons un $y_\lambda\in E_\lambda$ et
posons $y_\beta=g_{\beta\lambda}(y_\lambda)$,
$y_\alpha=g_{\alpha\lambda}(y_\lambda)$~; soit
$y'_\alpha\in g_{\alpha\beta}(E_\beta)$, de sorte que
$y'_\alpha=g_{\alpha\beta}(y'_\beta)$ pour un $y'_\beta\in E_\beta$~;
on a $y'_\beta-y_\beta=x_\beta\in A_\beta$, et par hypothèse il existe
$x_\lambda\in A_\lambda$ tel que
$g_{\alpha\beta}(x_\beta)=g_{\alpha\lambda}(x_\lambda)$~; donc
\[
y'_\alpha=g_{\alpha\beta}(y_\beta)+g_{\alpha\beta}(x_\beta)
=g_{\alpha\lambda}(y_\lambda)+g_{\alpha\lambda}(x_\lambda)
=g_{\alpha\lambda}(y_\lambda+x_\lambda)
\in g_{\alpha\lambda}(E_\lambda),
\]
ce qui démontre notre assertion. Cela étant, on sait (Bourbaki,
\emph{Top. gén.}, chap.~II, 3\textsuperscript{e} éd., §~3, th.~1) que sous
les hypothèses faites sur $I$, un système projectif d'ensembles non vides
vérifiant (ML) a une limite projective non vide~; par suite, il existe un
point $y=(y_\alpha)\in\varprojlim E_\alpha$, et comme
$v_\alpha(y_\alpha)=z_\alpha$ par définition pour tout $\alpha$, on a
$z=v(y)$, C.Q.F.D.
\end{proof}

\begin{proposition}[13.2.3]
\label{0.13.2.3-fr}
\emph{Les hypothèses sur $I$ étant celles de
(\hyperref[0.13.2.2-fr]{13.2.2}), soit
$(K_\alpha^\bullet)_{\alpha\in I}$ un système projectif de complexes de
groupes abéliens
$K_\alpha^\bullet=(K_\alpha^n)_{n\in\mathbb Z}$ dont l'opérateur de
dérivation est de degré $+1$. Pour chaque $n$, il existe un homomorphisme
canonique fonctoriel}
\begin{equation}
\label{0.13.2.3.1-fr}
\tag{13.2.3.1}
h_n:H^n(\varprojlim K_\alpha^\bullet)
\longrightarrow\varprojlim H^n(K_\alpha^\bullet).
\end{equation}
\emph{Si, pour tout degré $n$, le système projectif de groupes abéliens
$(K_\alpha^n)_{\alpha\in I}$ vérifie (ML), alors tous les homomorphismes
$h_n$ sont surjectifs. Si en outre, pour un degré $n$, le système
projectif $(H^{n-1}(K_\alpha^\bullet))_{\alpha\in I}$ vérifie (ML),
l'homomorphisme $h_n$ est bijectif.}
\end{proposition}

\begin{proof}
Posons, pour tout $n$, $K^n=\varprojlim K_\alpha^n$~; la définition des
homomorphismes $h_n$ provient de la commutativité des diagrammes
\[
\xymatrix{
\cdots\ar[r]&K^{n-1}\ar[r]\ar[d]&K^n\ar[r]\ar[d]
&K^{n+1}\ar[r]\ar[d]&\cdots\\
\cdots\ar[r]&K_\alpha^{n-1}\ar[r]&K_\alpha^n\ar[r]
&K_\alpha^{n+1}\ar[r]&\cdots
}
\]
